%
% reelle-beispiele.tex -- Reelle Beispiele und Gegenbeispiele
%
% !TEX root = ../../buch.tex
% !TEX encoding = UTF-8
%
\section{Reelle Beispiele und Gegenbeispiele
\label{hauptwert:section:reelle-beispiele}}
\kopfrechts{Reelle Beispiele}

\subsection{Ein weiteres Beispiel in der reellen Analysis}
Als zweites Beispiel betrachten wir die Funktion $f(x) = \frac{x + 1}{1 + x^2}$
und deren uneigentliches Integral $\int_{-\infty}^{\infty} \frac{x + 1}{1 + x^2}\,dx$.
\input{papers/hauptwert/fig/nichtungerade.tex}
Abbildung \ref{hauptwert:figure:nichtungerade} zeigt den Graphen des Integranden.
Versucht man wieder asymmetrisch von $-\infty$ nach $\infty$ zu integrieren,
erhält man
\begin{equation}
\begin{aligned}
  \int_{-\infty}^{\infty} \frac{x + 1}{1 + x^2}\,dx
  &= \int_{-\infty}^{\infty}
     \biggl(\frac{x}{1 + x^2} + \frac{1}{1 + x^2}\biggr)\,dx \\
  &= \Bigl[\tfrac{1}{2}\ln(1 + x^2) + \arctan x\Bigr]_{-\infty}^{\infty} \\
  &= \lim_{x\to\infty} \Bigl(\tfrac{1}{2}\ln(1 + x^2) + \arctan x\Bigr)
   - \lim_{x\to-\infty} \Bigl(\tfrac{1}{2}\ln(1 + x^2) + \arctan x\Bigr),
\end{aligned}
\end{equation}
also wieder den unbestimmten Ausdruck $\infty - \infty$.
Das uneigentliche Integral existiert also nicht.
Da der Integrand aber wieder ungerade ist, liefert die symmetrische
Annäherung nach Cauchy den Hauptwert
\begin{equation}
\begin{aligned}
  \pv \int_{-\infty}^{\infty} \frac{x + 1}{1 + x^2}\,dx
  &= \lim_{R\to\infty}
     \Bigl[\tfrac{1}{2}\ln(1 + x^2) + \arctan x\Bigr]_{-R}^{R} \\
  &= \lim_{R\to\infty} 2\arctan R \\
  &= \pi.
\end{aligned}
\end{equation}

\subsection{Eine Polstelle im Innern des Intervalls
\label{hauptwert:subsection:pol-innen}}
\input{papers/hauptwert/fig/pol-innen.tex}
Abbildung \ref{hauptwert:figure:pol-innen} zeigt den Graphen des Integranden.
\begin{equation}
\begin{aligned}
\pv\int_0^3 \frac{dx}{x-1}
  &= \lim_{\epsilon \to 0^+}
     \biggl(
     \int_{0}^{1 - \epsilon} \frac{dx}{x - 1} + \int_{1 + \epsilon}^{3} \frac{dx}{x - 1}
     \biggr) \\
  &= \lim_{\epsilon \to 0^+}
     \Bigl(\bigl[\ln|x - 1|\bigr]_{0}^{1 - \epsilon} + \bigl[\ln|x - 1|\bigr]_{1 + \epsilon}^{3}\Bigr) \\
  &= \lim_{\epsilon \to 0^+} \bigl(\ln\epsilon - \ln 1 + \ln 2 - \ln\epsilon\bigr) \\
  &= \ln 2
\label{hauptwert:equation:beispiel-pol-innen}
\end{aligned}
\end{equation}
